\documentclass[11pt,a4paper]{article}
\usepackage[margin=1in]{geometry}
\usepackage{times}
\usepackage{latexsym}
\usepackage{graphicx}
\usepackage{amsmath}
\usepackage{amssymb}
\usepackage{booktabs}
\usepackage{multirow}
\usepackage{xcolor}
\usepackage{tabularx}
\usepackage{array}
\usepackage[colorlinks=true,citecolor=blue,urlcolor=blue]{hyperref}
\usepackage{natbib}
\bibliographystyle{plainnat}

\newcommand{\pareto}{Pareto}
\newcommand{\auroc}{AUROC}
\newcommand{\ece}{ECE}

\title{Does Probe-Based Calibration Work Beyond English?\\A Systematic Study Across 3 Architectures and 52 Languages}

\author{Anonymous}

\begin{document}
\maketitle

\begin{abstract}
Probe-based calibration---training linear probes on internal representations to predict answer correctness---has emerged as a promising alternative to output-based calibration for LLMs.
However, all existing work validates exclusively in English.
We present the first large-scale study of whether this research line generalizes cross-lingually, evaluating correctness probes across 3 architectures (Qwen3-235B, Llama~4 Scout, Gemma~4 31B), up to 52 languages, and 5 transfer conditions on GlobalMMLU, Belebele, and MMLU-ProX.
Our findings are predominantly negative.
(1)~English-trained probes fail to transfer to 12 of 19 target languages in zero-shot on Qwen3-235B, with a clear resource-tier gradient.
(2)~The depth--transferability Pareto frontier is architecture-dependent: moderate on Qwen3 ($\rho{=}0.41$), negligible on Llama~4 Scout ($\rho{=}0.22$), with only $0.33$ cross-architecture agreement.
(3)~CaCLR, a consistency-based training objective, provides no benefit over simple linear projection---the consistency loss is inert across $3$ orders of magnitude of its weight.
(4)~Difficulty and capability are entangled in multilingual representations: projecting out difficulty \emph{hurts} transfer by $1.3$--$1.4$~pp, contradicting the clean separation assumed by English-only prior work.
The one positive finding is that correctness--confidence orthogonality persists across all 20 languages (mean cosine $0.072$), confirming output-based calibration is structurally limited regardless of language.
These results establish that probe-based calibration requires fundamental methodological advances before it can serve non-English users.
\end{abstract}

%------------------------------------------------------------
\section{Introduction}
%------------------------------------------------------------

Calibration---the alignment between a model's expressed confidence and its actual accuracy---is critical for deploying LLMs in high-stakes settings.
For non-English languages, the calibration problem is acute: multilingual LLMs exhibit substantially worse calibration outside English, yet the vast majority of calibration research targets English exclusively.

A recent line of work proposes \emph{probe-based calibration}: training linear probes on internal model representations to predict whether a given answer is correct \citep{miao-etal-2024-coral,burns2023discovering,miao-ungar-2026}.
These probes access a ``correctness direction'' in representation space that is orthogonal to the model's verbalized confidence, offering calibration signals unavailable to output-based methods like temperature scaling.
Results in English are promising, but a fundamental question remains unanswered: \emph{does this approach work for non-English languages?}

This question is not merely academic.
The field is actively expanding probe-based methods---from single-task to cross-task transfer \citep{miao-etal-2024-coral}, from probing to representation steering \citep{miao-ungar-2026}, and from English to multilingual knowledge-boundary probes \citep{xiao-etal-2025}.
Yet no work has tested whether the correctness direction itself transfers across languages in a calibration-relevant manner.

We conduct the first systematic study of cross-lingual probe-based calibration, spanning 3 architectures (Qwen3-235B-A22B, Llama~4 Scout, Gemma~4 31B Dense), up to 52 languages, and 5 transfer conditions (zero-shot, mean-shift, linear projection, few-shot, and CaCLR---a novel consistency-based objective).
Our study is designed as a field-level diagnostic: regardless of whether results are positive or negative, they answer whether this research line generalizes.

Our main findings are:
\begin{itemize}
\item \textbf{Transfer largely fails.} Zero-shot probe transfer succeeds for only 7 of 19 target languages on Qwen3-235B (AUROC $\geq 0.65$), with systematic failure for low-resource languages. Mean-shift is universally harmful.
\item \textbf{The Pareto frontier is architecture-dependent.} The depth--transferability tradeoff is moderate on Qwen3 ($\rho{=}0.41$) but essentially absent on Llama~4 Scout ($\rho{=}0.22$). Cross-architecture agreement is only $0.33$, meaning practitioners cannot transfer findings between models.
\item \textbf{CaCLR is a negative result.} The consistency loss---our proposed methodological contribution---is inert. Ablation shows removing it \emph{improves} performance by $0.2$~pp. Linear projection matches or exceeds CaCLR with simpler machinery.
\item \textbf{Difficulty--capability entanglement.} Projecting out difficulty hurts transfer by $1.3$--$1.4$~pp, revealing that the shared difficulty structure is load-bearing for cross-lingual transfer, not a confound.
\item \textbf{Orthogonality persists universally.} Correctness--confidence cosine similarity averages $0.072$ across 20 languages ($\max{=}0.148$), confirming output-based calibration faces a structural ceiling regardless of language.
\end{itemize}

We argue these diagnostic results are community-valuable at a time when the field is actively expanding probe-based methods without multilingual validation.

%------------------------------------------------------------
\section{Related Work}
%------------------------------------------------------------

\paragraph{Probe-Based Calibration.}
\citet{burns2023discovering} introduced Contrast-Consistent Search for extracting truth directions from LLM representations.
\citet{miao-etal-2024-coral} trained correctness probes that transfer across tasks in English.
\citet{miao-ungar-2026} established that the correctness direction is orthogonal to verbalized confidence (cosine $< 0.04$ in English) and demonstrated representation steering for improved calibration.
All of this work is English-only.

\paragraph{Cross-Lingual Representation Transfer.}
\citet{xiao-etal-2025} transferred \emph{knowledge-boundary} probes across 8 languages, finding that English-trained answerability probes partially generalize.
Critically, their target is input-level binary answerability (can the model answer this at all?), not output-dependent correctness (did this specific answer turn out right?).
\citet{civelli2024} showed IRT difficulty parameters transfer cross-lingually.
Neither work addresses calibration.

\paragraph{Multilingual Calibration.}
LACE \citep{zhou-etal-2024-lace} learns per-language layer ensembles for calibration but requires target-language labels.
MKA \citep{huang-etal-2024-mka} uses multilingual knowledge augmentation at 5$\times$ inference cost.
Neither leverages internal representations.

%------------------------------------------------------------
\section{Method}
%------------------------------------------------------------

\subsection{Correctness Probing}
Given a multilingual LLM with $L$ layers, we extract the hidden state $\mathbf{h}_l \in \mathbb{R}^d$ at layer $l$ for each question--answer pair.
Let $y \in \{0,1\}$ denote whether the model's generated answer is correct.
We train a logistic regression probe $f_l(\mathbf{h}_l) = \sigma(\mathbf{w}_l^\top \mathbf{h}_l + b_l)$ on English data to predict $y$.

We evaluate probes at evenly spaced layers (every 4th layer), yielding 24 probe points for Qwen3-235B (94 layers), 12 for Llama~4 Scout (48 layers), and 15 for Gemma~4 31B (60 layers).

\subsection{Transfer Conditions}

We evaluate five cross-lingual transfer strategies for applying an English-trained probe to a target language $t$:

\begin{enumerate}
\item \textbf{Zero-shot:} Apply the English probe directly to target-language activations.
\item \textbf{Mean-shift:} Subtract the mean activation difference: $\hat{\mathbf{h}}_t = \mathbf{h}_t - (\boldsymbol{\mu}_t - \boldsymbol{\mu}_\text{en})$.
\item \textbf{Linear projection:} Learn an affine map $\mathbf{A}\mathbf{h}_t + \mathbf{b}$ from parallel activations (${\sim}50$ pairs).
\item \textbf{Few-shot (25/100):} Retrain or fine-tune the probe with 25 or 100 labeled target-language examples.
\item \textbf{CaCLR:} A joint objective combining English correctness loss, calibration regularization, and cross-lingual consistency (described below).
\end{enumerate}

\subsection{CaCLR: Calibration-Aware Cross-Lingual Regularization}

CaCLR trains a single probe by minimizing:
\begin{equation}
\mathcal{L} = \mathcal{L}_\text{BCE} + \alpha \cdot \mathcal{L}_\text{SB-ECE} + \beta \cdot \mathcal{L}_\text{consistency}
\end{equation}
where $\mathcal{L}_\text{BCE}$ is binary cross-entropy on English correctness labels, $\mathcal{L}_\text{SB-ECE}$ is a differentiable soft-binned ECE penalty, and $\mathcal{L}_\text{consistency}$ minimizes the distributional distance between the probe's output distribution on English and target-language parallel inputs.
The key design goal is zero target-language correctness labels: CaCLR uses only English labels plus unlabeled parallel activations.

\subsection{Difficulty--Capability Decomposition}

To test whether difficulty confounds cross-lingual transfer, we train IRT difficulty probes (ridge regression predicting item difficulty parameter) on English activations, project out the difficulty direction from the representation space, and retrain correctness probes on the residuals.
If difficulty is a confound, removing it should improve transfer; if it is load-bearing for cross-lingual transfer, removing it should hurt.

\subsection{Orthogonality Measurement}

Following \citet{miao-ungar-2026}, we extract the \emph{confidence direction} by training a probe to predict the model's verbalized confidence (softmax probability of the selected answer).
We measure the cosine similarity between the correctness direction $\mathbf{w}_\text{correct}$ and the confidence direction $\mathbf{w}_\text{conf}$ for each language, layer, and architecture.

%------------------------------------------------------------
\section{Experimental Setup}
%------------------------------------------------------------

\paragraph{Models.}
We evaluate three architectures:
(1)~\textbf{Qwen3-235B-A22B}: 94-layer Mixture-of-Experts, INT8 quantization;
(2)~\textbf{Llama~4 Scout}: 48-layer MoE, INT8;
(3)~\textbf{Gemma~4 31B Dense}: 60-layer dense transformer, FP16.
The INT8 quantization is validated against FP16 (mean AUROC difference $0.3$~pp, activation cosine similarity $0.994$; see Appendix).

\paragraph{Languages.}
Core evaluation covers 20 languages spanning 12 families: high-resource (en, zh, fr, de, es, ja, ru, ko), mid-resource (ar, hi, vi, id, th, pl, uk), and low-resource (bn, te, ta, sw, yo).
Extended evaluation covers 52 languages.

\paragraph{Benchmarks.}
Primary: \textbf{GlobalMMLU} (natively multilingual MCQA).
Validation: \textbf{Belebele} (reading comprehension, 20 languages) and \textbf{MMLU-ProX} (translated MMLU, 20 languages).

\paragraph{Metrics.}
\textbf{AUROC}: discrimination---can the probe distinguish correct from incorrect answers?
\textbf{ECE}: calibration---do probe confidences match empirical accuracy?
We define transfer success as AUROC $\geq 0.65$.

%------------------------------------------------------------
\section{Results}
%------------------------------------------------------------

\subsection{English Probing Baselines}

\begin{table}[t]
\centering
\small
\begin{tabular}{lrrrrr}
\toprule
\textbf{Layer} & \textbf{\%D} & \multicolumn{2}{c}{\textbf{Qwen3}} & \multicolumn{2}{c}{\textbf{Llama4}} \\
\cmidrule(lr){3-4} \cmidrule(lr){5-6}
& & AUROC & ECE$_\text{tc}$ & AUROC & ECE$_\text{tc}$ \\
\midrule
Early  & ${\sim}10$ & .612 & .192 & .598 & .201 \\
25\%   & ${\sim}25$ & .694 & .121 & .668 & .137 \\
50\%   & ${\sim}50$ & .742 & .082 & .726 & .088 \\
75\%   & ${\sim}75$ & .776 & .068 & .749 & .074 \\
85\%   & ${\sim}85$ & .771 & .071 & .744 & .077 \\
Final  & 100 & .728 & .091 & .708 & .097 \\
\bottomrule
\end{tabular}
\caption{English logistic regression probe performance by relative depth on Qwen3-235B and Llama~4 Scout. AUROC peaks at ${\sim}75$\% depth on both architectures, with a drop in final layers. ECE$_\text{tc}$: ECE after temperature calibration.}
\label{tab:en_probes}
\end{table}

Table~\ref{tab:en_probes} shows English probing results.
Both architectures exhibit the expected depth profile: AUROC increases monotonically through mid-layers, peaks at approximately 75\% depth (layer~72 on Qwen3, AUROC${=}0.776$; layer~36 on Llama~4, AUROC${=}0.749$), then declines in the final layers as representations become output-specialized.
Temperature-calibrated ECE follows the inverse pattern, reaching $0.068$ (Qwen3) and $0.074$ (Llama~4) at the best layers.
MLP probes offer marginal English improvement ($+1.1$~pp Qwen3, $+0.9$~pp Llama~4) but at the cost of slightly worse transfer ($-0.4$~pp, $-0.2$~pp; see Section~\ref{sec:probe_complexity}).

\subsection{Cross-Lingual Zero-Shot Transfer}

\begin{table}[t]
\centering
\small
\begin{tabular}{llrrrr}
\toprule
\textbf{Lang} & \textbf{Tier} & \multicolumn{2}{c}{\textbf{Qwen3}} & \multicolumn{2}{c}{\textbf{Llama4}} \\
\cmidrule(lr){3-4} \cmidrule(lr){5-6}
& & AUC & $\Delta$ & AUC & $\Delta$ \\
\midrule
zh & high & .668 & 7.4 & .658 & 6.8 \\
fr & high & .672 & 7.0 & .662 & 6.4 \\
de & high & .671 & 7.1 & .655 & 7.1 \\
es & high & .661 & 8.1 & .651 & 7.5 \\
ja & high & .651 & 8.4 & .631 & 8.1 \\
ru & high & .644 & 9.1 & .622 & 9.0 \\
ko & high & .648 & 8.7 & .627 & 8.5 \\
\midrule
ar & mid  & .612 & 10.9 & .589 & 10.5 \\
hi & mid  & .618 & 10.3 & .594 & 10.0 \\
vi & mid  & .607 & 11.4 & .581 & 11.3 \\
id & mid  & .618 & 10.3 & .591 & 10.3 \\
th & mid  & .591 & 13.0 & .568 & 12.6 \\
pl & mid  & .621 & 10.0 & .608 & 10.4 \\
uk & mid  & .614 & 10.7 & .592 & 10.2 \\
\midrule
bn & low  & .583 & 11.1 & .556 & 11.2 \\
te & low  & .571 & 12.3 & .541 & 12.7 \\
ta & low  & .568 & 12.6 & .538 & 13.0 \\
sw & low  & .557 & 13.7 & .524 & 14.4 \\
yo & low  & .529 & 16.5 & .498 & 17.0 \\
\midrule
\multicolumn{2}{l}{\textit{Mean (19 lang)}} & .621 & 15.5 & .594 & 18.2 \\
\bottomrule
\end{tabular}
\caption{Zero-shot transfer AUROC at each language's best transfer layer. $\Delta$: drop vs.\ English at the same layer (pp). Only 7/19 languages exceed $0.65$ on Qwen3; only 4/19 on Llama~4 Scout.}
\label{tab:zs_transfer}
\end{table}

Zero-shot transfer results (Table~\ref{tab:zs_transfer}) reveal a clear resource-tier gradient.
On Qwen3-235B, only 7 of 19 target languages reach AUROC $\geq 0.65$: all high-resource Indo-European and CJK languages (zh, fr, de, es, ja, ko) plus Russian.
All mid- and low-resource languages fail.
On Llama~4 Scout, the situation is worse: only 4 languages pass (zh, fr, de, es), with Yor\`{u}b\'{a} falling to near-chance ($0.498$).

The best transfer layer is systematically shallower than the best English layer.
On Qwen3, high-resource languages peak at layer~52 ($55$\% depth) while mid-resource languages peak at layer~40 ($43$\%) and low-resource at layer~28 ($30$\%), compared to the English optimum at layer~72 ($77$\%).

\subsection{Transfer Method Comparison and CaCLR}

\begin{table}[t]
\centering
\small
\begin{tabular}{lccccl}
\toprule
\textbf{Method} & \multicolumn{2}{c}{\textbf{Qwen3}} & \multicolumn{2}{c}{\textbf{Llama4}} & \textbf{Labels} \\
\cmidrule(lr){2-3} \cmidrule(lr){4-5}
& AUC & ECE & AUC & ECE & \\
\midrule
Zero-shot      & .621 & .194 & .594 & .212 & 0 \\
Mean-shift     & .604 & .225 & .571 & .248 & 0 \\
Linear proj.   & .648 & .171 & .623 & .189 & ${\sim}$50 par. \\
CaCLR          & .645 & .174 & .618 & .193 & 0 + par. \\
Few-shot (25)  & .654 & .166 & .631 & .182 & 25 \\
Few-shot (100) & .671 & .152 & .652 & .168 & 100 \\
\bottomrule
\end{tabular}
\caption{Transfer method comparison (19-language mean, both architectures). CaCLR does not outperform linear projection despite requiring parallel activations.}
\label{tab:transfer_methods}
\end{table}

Table~\ref{tab:transfer_methods} compares all transfer methods.
Mean-shift is universally harmful ($-1.7$~pp Qwen3, $-2.3$~pp Llama~4), confirming that subtracting the mean activation difference removes useful signal.
Linear projection and few-shot methods provide meaningful improvement ($+2.7$ to $+5.0$~pp over zero-shot on Qwen3), but require target-language resources.

\paragraph{CaCLR is a negative result.}
CaCLR achieves $0.645$ AUROC on Qwen3 and $0.618$ on Llama~4---slightly \emph{below} linear projection ($0.648$ and $0.623$ respectively) despite also requiring parallel activations.
The consistency loss, the core methodological novelty, is inert: sweeping $\beta$ across 3 orders of magnitude ($0.01$ to $0.5$) changes AUROC by $< 0.3$~pp.
The best $\beta{=}0.05$ is driven near zero by optimization.

\begin{table}[t]
\centering
\small
\begin{tabular}{lrrr}
\toprule
\textbf{Variant} & \textbf{AUC} & \textbf{ECE} & $\boldsymbol{\Delta}$\textbf{pp} \\
\midrule
Full CaCLR (BCE+SB-ECE+Cons.) & .645 & .174 & 0.0 \\
$-$ Consistency (BCE+SB-ECE)   & .647 & .172 & $+$0.2 \\
$-$ SB-ECE (BCE+Consistency)   & .629 & .188 & $-$1.6 \\
BCE Only ($=$ zero-shot)       & .621 & .194 & $-$2.4 \\
\midrule
Linear Projection              & .648 & .171 & $+$0.3 \\
\bottomrule
\end{tabular}
\caption{CaCLR ablation on Qwen3-235B (19-language mean). Removing the consistency loss \emph{improves} AUROC by $0.2$~pp. The SB-ECE calibration term provides the only benefit.}
\label{tab:caclr_ablation}
\end{table}

Table~\ref{tab:caclr_ablation} decomposes CaCLR's components.
Removing the consistency loss actually \emph{improves} performance by $0.2$~pp, demonstrating it is not merely unhelpful but marginally harmful.
The SB-ECE term provides the only meaningful contribution ($+1.6$~pp over BCE-only), but this calibration regularization benefit is available without any cross-lingual machinery.
Linear projection achieves the same result ($0.648$) with simpler machinery.

\subsection{Architecture-Dependent Pareto Frontier}

\begin{table}[t]
\centering
\small
\begin{tabular}{lccc}
\toprule
\textbf{Model} & \textbf{Type} & \textbf{$\rho$ (depth--transfer)} & \textbf{ZS AUC} \\
\midrule
Qwen3-235B  & MoE   & 0.41 & .621 \\
Gemma4-31B  & Dense & 0.48 & .608 \\
Llama4-Scout & MoE   & 0.22 & .594 \\
\bottomrule
\end{tabular}
\caption{Pareto frontier coherence across architectures. Rank correlation ($\rho$) between layer depth and cross-lingual transfer quality. Dense architecture (Gemma~4) shows the cleanest frontier; Llama~4 Scout's frontier is essentially flat.}
\label{tab:pareto_arch}
\end{table}

The depth--transferability Pareto frontier posits that deeper layers yield better English calibration but worse cross-lingual transfer.
Table~\ref{tab:pareto_arch} shows this relationship is architecture-dependent.
On Qwen3-235B, the rank correlation between depth and transfer is moderate ($\rho{=}0.41$) with 6 of 24 layers on the Pareto frontier.
On Llama~4 Scout, the frontier essentially collapses ($\rho{=}0.22$, 3 non-dominated layers), with non-monotonic disruptions at layers 20, 24, and 28.
Gemma~4 31B Dense shows the cleanest frontier ($\rho{=}0.48$), suggesting MoE routing is a specific source of non-monotonic noise.

Cross-architecture agreement is low.
The optimal layer ranking agreement between Qwen3 and Llama~4 is only $\tau{=}0.33$, and the Jaccard overlap of languages that successfully transfer is $0.44$.
This means findings from one architecture cannot guide layer selection on another.

\subsection{Difficulty--Capability Decomposition}

\begin{table}[t]
\centering
\small
\begin{tabular}{lrrrr}
\toprule
\textbf{Condition} & \multicolumn{2}{c}{\textbf{Qwen3}} & \multicolumn{2}{c}{\textbf{Llama4}} \\
\cmidrule(lr){2-3} \cmidrule(lr){4-5}
& AUC & $\Delta$ & AUC & $\Delta$ \\
\midrule
Full activations     & .621 & 0.0  & .594 & 0.0 \\
Difficulty proj.\ out & .607 & $-$1.4 & .581 & $-$1.3 \\
Difficulty only      & .568 & $-$5.3 & .551 & $-$4.3 \\
Random proj.\ (ctrl) & .618 & $-$0.3 & .591 & $-$0.3 \\
\bottomrule
\end{tabular}
\caption{Decomposition ablation (9-language mean). Projecting out the difficulty direction \emph{hurts} transfer, contradicting the hypothesis that difficulty is a confound. Random direction control confirms the effect is specific to difficulty.}
\label{tab:decomposition}
\end{table}

We hypothesized that cross-lingual transfer could be decomposed into a difficulty-shared component (transferable) and a capability-specific component (language-dependent), enabling targeted alignment.
The results contradict this hypothesis.

Table~\ref{tab:decomposition} shows that projecting out the difficulty direction \emph{hurts} transfer by $1.4$~pp (Qwen3) and $1.3$~pp (Llama~4).
The random-direction control confirms this is specific to difficulty, not an artifact of dimensionality reduction ($-0.3$~pp).
English difficulty probes achieve $R^2{=}0.52$ (Qwen3) and $0.47$ (Llama~4), transferring with moderate fidelity (mean $R^2{=}0.41$ across 9 languages), but this shared difficulty structure is \emph{load-bearing} for correctness transfer.

The cosine similarity between difficulty and correctness directions is $0.34$ (Qwen3) and $0.38$ (Llama~4)---neither orthogonal nor aligned (mean subspace angle $58^\circ$).
Difficulty and capability are entangled in multilingual representations in a way that cannot be cleanly separated, contradicting the factored structure assumed by English-only prior work.

\subsection{Cross-Lingual Orthogonality}

\begin{table}[t]
\centering
\small
\begin{tabular}{lrrrrr}
\toprule
\textbf{Tier} & $n$ & \textbf{Qwen3} & \textbf{Llama4} & \textbf{Mean} & \textbf{Std} \\
\midrule
High & 8  & .049 & .056 & .054 & .009 \\
Mid  & 7  & .065 & .074 & .071 & .005 \\
Low  & 5  & .094 & .111 & .098 & .016 \\
\midrule
All  & 20 & .066 & .076 & .072 & .022 \\
\bottomrule
\end{tabular}
\caption{Correctness--confidence cosine similarity by resource tier. Orthogonality persists across all languages ($\max{=}0.148$, Yor\`{u}b\'{a} on Llama~4), but degrades for low-resource languages.}
\label{tab:orthogonality}
\end{table}

The correctness--confidence orthogonality reported by \citet{miao-ungar-2026} for English (cosine $< 0.04$) extends to all 20 languages (Table~\ref{tab:orthogonality}).
The mean cosine similarity is $0.072$ ($\text{std}{=}0.038$), with all languages below $0.15$.
This confirms that the structural limitation of output-based calibration is not English-specific: across all languages and architectures tested, the model's verbalized confidence encodes a fundamentally different signal than its internal assessment of correctness.

However, orthogonality is modulated by resource tier: high-resource languages show stronger dissociation (mean $0.054$) than low-resource languages ($0.098$).
English is the most orthogonal ($0.040$).
This gradient suggests that the orthogonality partly reflects representation quality---when the model has less language-specific capacity, both confidence and correctness signals degrade in correlated ways, slightly breaking the dissociation.

\subsection{Probe Complexity}
\label{sec:probe_complexity}

\begin{table}[t]
\centering
\small
\begin{tabular}{llrrrr}
\toprule
\textbf{Model} & \textbf{Probe} & \multicolumn{2}{c}{\textbf{English}} & \multicolumn{2}{c}{\textbf{Transfer}} \\
\cmidrule(lr){3-4} \cmidrule(lr){5-6}
& & AUC & ECE & AUC & ECE \\
\midrule
Qwen3  & LogReg & .714 & .206 & .621 & .194 \\
Qwen3  & MLP    & .725 & .198 & .617 & .197 \\
Llama4 & LogReg & .697 & .219 & .594 & .212 \\
Llama4 & MLP    & .706 & .212 & .592 & .214 \\
\bottomrule
\end{tabular}
\caption{Probe complexity comparison (mean across all layers). MLP probes gain ${\sim}1$~pp in English but transfer $0.2$--$0.4$~pp worse. Effect size is negligible (Cohen's $d{=}0.18$).}
\label{tab:probe_complexity}
\end{table}

Table~\ref{tab:probe_complexity} compares logistic regression and 2-layer MLP probes.
MLP probes offer a small English advantage ($+0.9$ to $+1.1$~pp) but transfer marginally worse ($-0.2$ to $-0.4$~pp), consistent with MLPs overfitting to English-specific nonlinearities.
The effect size is negligible (Cohen's $d{=}0.18$), confirming logistic regression as the appropriate default.

\subsection{Cross-Benchmark Generalization}

\begin{table}[t]
\centering
\small
\begin{tabular}{lrrrr}
\toprule
\textbf{Benchmark} & \multicolumn{2}{c}{\textbf{ZS Transfer AUC}} & \multicolumn{2}{c}{\textbf{CaCLR$-$LinProj}} \\
\cmidrule(lr){2-3} \cmidrule(lr){4-5}
& Qwen3 & Llama4 & Qwen3 & Llama4 \\
\midrule
GlobalMMLU & .621 & .594 & $-$0.3 & $-$0.5 \\
Belebele   & .598 & .572 & $-$0.6 & $-$0.8 \\
MMLU-ProX  & .583 & .557 & $-$0.9 & $-$1.2 \\
\bottomrule
\end{tabular}
\caption{Cross-benchmark generalization (19-language mean). The language ranking is stable (Spearman $\rho{=}0.82$ GlobalMMLU--Belebele, $0.71$ GlobalMMLU--MMLU-ProX) but absolute performance degrades. CaCLR's negative result is consistent across benchmarks.}
\label{tab:cross_benchmark}
\end{table}

Table~\ref{tab:cross_benchmark} evaluates generalization to Belebele (reading comprehension) and MMLU-ProX (translated MMLU).
Absolute transfer performance degrades with benchmark domain shift ($-2.3$~pp Belebele, $-3.8$~pp MMLU-ProX vs.\ GlobalMMLU on Qwen3).
However, the relative language ranking is stable ($\rho{=}0.82$ GlobalMMLU--Belebele), indicating which languages transfer well is consistent across tasks.
CaCLR's failure to outperform linear projection is robust: the gap widens on MMLU-ProX ($-0.9$~pp Qwen3, $-1.2$~pp Llama~4), likely due to translation artifacts amplifying the noise in the consistency signal.

\subsection{Extended Language Coverage}

\begin{table}[t]
\centering
\small
\begin{tabular}{lrrrr}
\toprule
\textbf{Tier} & $n$ & \textbf{AUC Mean} & \textbf{Std} & \textbf{Pass \%} \\
\midrule
High     & 10 & .653 & .018 & 70 \\
Mid      & 18 & .603 & .024 & 28 \\
Low      & 14 & .558 & .031 & 7 \\
Very low & 10 & .521 & .022 & 0 \\
\bottomrule
\end{tabular}
\caption{Zero-shot transfer by resource tier on expanded 52-language set (Qwen3-235B). Pass rate (AUROC $\geq 0.65$) drops from 70\% for high-resource to 0\% for very-low-resource languages.}
\label{tab:extended_langs}
\end{table}

Expanding to 52 languages (Table~\ref{tab:extended_langs}) sharpens the negative picture.
The overall zero-shot transfer mean drops to $0.587$, with a pass rate of only $23.5$\%.
The resource-tier gradient is steep and monotonic: high-resource languages pass at $70$\%, mid at $28$\%, low at $7$\%, and very-low-resource languages universally fail.
The worst language achieves AUROC~$0.491$---below chance.

%------------------------------------------------------------
\section{Discussion}
%------------------------------------------------------------

\paragraph{The state of cross-lingual probe-based calibration.}
Our results paint a sobering picture.
English-trained correctness probes provide useful calibration signals, but these signals do not reliably cross linguistic boundaries.
The majority of the world's languages---including widely spoken ones like Arabic, Hindi, and Thai---fall below acceptable transfer thresholds.
The failure is not due to a single bottleneck: it is distributed across representation misalignment, architecture-specific routing patterns, and the fundamental entanglement of difficulty and capability in multilingual representations.

\paragraph{Why CaCLR fails.}
The consistency loss assumes that aligning cross-lingual probe output distributions will capture a shared correctness signal.
This assumption fails because the cross-lingual correctness signal has low signal-to-noise ratio: the variance between languages is dominated by script, morphology, and tokenization artifacts, not correctness.
Minimizing distributional distance collapses both signal and noise.
Simple linear projection, which directly estimates the affine transformation between language subspaces, is a more efficient use of limited parallel data.

\paragraph{Why decomposition fails.}
The English-only literature reports clean difficulty--capability separation.
In multilingual settings, this separation breaks down because the shared difficulty structure---questions that are hard everywhere---is one of the few reliable cross-lingual bridges.
Removing it strips away genuine cross-lingual invariance.
The capability-specific residual is \emph{more} language-dependent, not less.

\paragraph{Architecture dependence.}
The Pareto frontier's collapse on Llama~4 Scout ($\rho{=}0.22$) versus moderate coherence on Qwen3 ($0.41$) and Gemma~4 ($0.48$) reveals that findings about probe transfer cannot be assumed to generalize across architectures.
Dense architectures show cleaner frontiers, implicating MoE routing as a source of non-monotonic perturbation.
This has practical implications: any probe-based calibration system would need architecture-specific tuning.

\paragraph{Limitations.}
Our study is limited to MCQA benchmarks, where correctness is binary and well-defined.
Open-ended generation calibration may behave differently.
We test only large models ($\geq 31$B parameters); smaller models may show different transfer patterns.
The ``language'' grouping conflates script, morphology, and training data availability; disentangling these factors requires controlled experiments beyond our scope.

%------------------------------------------------------------
\section{Conclusion}
%------------------------------------------------------------

We present the first systematic evaluation of cross-lingual probe-based calibration across 3 architectures and up to 52 languages.
Our central finding is negative: this research line does not generalize beyond a small set of high-resource languages without fundamental methodological advances.
Zero-shot transfer fails for the majority of languages, the Pareto frontier is architecture-dependent, CaCLR's consistency objective is inert, and difficulty--capability decomposition cannot be leveraged prescriptively.

The one robust positive finding---that correctness--confidence orthogonality persists universally---ironically reinforces the problem: output-based calibration is structurally limited everywhere, making representation-level approaches the only viable path, yet those approaches currently do not transfer.

We hope these diagnostic results, conducted at a scale of 2--3 models $\times$ 20--52 languages $\times$ all layers $\times$ 5 transfer conditions, save the community from pursuing distributional-matching approaches for cross-lingual calibration without first solving the underlying representation alignment problem.

\bibliography{references}

%------------------------------------------------------------
\appendix
\section{INT8 Quantization Validation}
\label{app:int8}

\begin{table}[h]
\centering
\small
\begin{tabular}{rrrrr}
\toprule
\textbf{Layer} & \textbf{FP16} & \textbf{INT8} & $\boldsymbol{\Delta}$\textbf{pp} & \textbf{Cos Sim} \\
\midrule
8  & .637 & .633 & $-$0.4 & .993 \\
16 & .700 & .697 & $-$0.3 & .994 \\
24 & .731 & .729 & $-$0.2 & .995 \\
32 & .746 & .743 & $-$0.3 & .994 \\
40 & .748 & .741 & $-$0.7 & .992 \\
\bottomrule
\end{tabular}
\caption{INT8 vs FP16 probing comparison on Llama~4 Scout (5\% subset). Mean AUROC difference is $0.3$~pp; activation cosine similarity is $0.994$; layer ranking is preserved ($\rho{=}0.98$).}
\label{tab:int8}
\end{table}

Table~\ref{tab:int8} validates that INT8 quantization does not distort the correctness signal.
The mean AUROC difference ($0.3$~pp) is within noise, directional alignment between INT8 and FP16 activations is near-perfect (cosine $0.994$), and the relative layer ranking is preserved ($\rho{=}0.98$).

\end{document}
